\subsection{Passage to the limit in integrals}

PROPOSITION 3. --- Let B be a filter base on \(L_{F}^{1}\) . Assume that there exists a compact set \(K \subset X\) such that, for every set \(M \in B\) , all of the functions \(f \in M\) have their support in K. Under these conditions, if B converges uniformly on X to \(f_{0}\) , then the function \(f_{0}\) is integrable and

\[
\int \mathbf {f} _ {0} d \mu = \lim _ {\mathfrak {B}} \int \mathbf {f} d \mu .\tag{9}
\]

For, \(\mathfrak{B}\) converges in mean to \(\mathbf{f}_0\) (§3, No.~3, Prop. 4).

PROPOSITION 4. --- Let \((f_n)\) be an increasing (resp. decreasing) sequence of integrable numerical functions. For the upper (resp. lower) envelope \(f\) of the sequence to be integrable, it is necessary and sufficient that \(\sup_{n} \int f_n d|\mu| < +\infty\) (resp. \(\inf_{n} \int f_n d|\mu| > -\infty\) ), in which case

\[
\int f d \mu = \lim _ {n \to \infty} \int f _ {n} d \mu .\tag{10}
\]

We limit ourselves to considering an increasing sequence. The sequence \(g_{n} = f_{n} + f_{1}^{-}\) is increasing and consists of integrable functions \(\geqslant 0\) ; since its upper envelope is \(g = f + f_{1}^{-}\) , the proposition follows from Th. 5 of §3, No.~6.

THEOREM 2. --- Let A be a set of indices, filtered by a filter F with a countable base. Let \((\mathbf{f}_{\alpha})_{\alpha\in\mathrm{A}}\) be a family of integrable functions that, with respect to the filter F, converge pointwise almost everywhere to a function f; if there exists a numerical function \(g\geqslant0\) such that \(\int^{*}gd|\mu|<+\infty\) and such that \(|\mathbf{f}_{\alpha}(x)|\leqslant g(x)\) almost everywhere in X for each \(\alpha\in A\) , then the function f is integrable and

\[
\int \mathbf {f} d \mu = \lim _ {\mathfrak {F}} \int \mathbf {f} _ {\alpha} d \mu .\tag{11}
\]

The theorem follows from Lebesgue's theorem (§3, No.~7, Cor. of Th. 6) since, under the conditions of the statement, \(\mathbf{f}_{\alpha}\) converges in mean to \(\mathbf{f}\) with respect to \(\mathfrak{F}\) .

COROLLARY 1. --- Let \(\Omega\) be a topological space, \(t_0\) a point of \(\Omega\) admitting a countable fundamental system of neighborhoods, \(\mathbf{f}\) a mapping of \(\mathbf{X} \times \Omega\) into \(\mathbf{F}\) having the following properties:

\begin{enumerate}
\def\labelenumi{\alph{enumi})}
\item
  for every \(t \in \Omega\) , the function \(x \mapsto \mathbf{f}(x, t)\) is integrable;
\item
  for every \(x \in \mathrm{X}\) , the function \(t \mapsto \mathbf{f}(x, t)\) is continuous at \(t_0\) ;
\item
  there exist a neighborhood U of \(t_0\) and a numerical function \(g \geqslant 0\) defined on X, such that \(\int^{*} g d|\mu| < +\infty\) and \(|\mathbf{f}(x,t)| \leqslant g(x)\) for all \(x \in X\) and \(t \in U\) .
\end{enumerate}

Under these conditions, the mapping \(t \mapsto \int \mathbf{f}(x,t) d\mu(x)\) of \(\Omega\) into \(\mathbf{F}\) is continuous at the point \(t_0\) .

COROLLARY 2. --- Let \((\mathbf{f}_{n})\) be a sequence of integrable functions such that the series with general term \(\mathbf{f}_{n}(x)\) converges almost everywhere; if there exists a function \(g \geqslant 0\) such that \(\int^{*} g d|\mu| < +\infty\) and such that, for every integer n, \(\left|\sum_{k=1}^{n} \mathbf{f}_{k}(x)\right| \leqslant g(x)\) almost everywhere, then the sum \(\mathbf{f}(x)\) (defined almost everywhere) of the series with general term \(\mathbf{f}_{n}(x)\) is integrable and

\[
\int \mathbf {f} d \mu = \sum_ {n = 1} ^ {\infty} \int \mathbf {f} _ {n} d \mu\tag{12}
\]

(`term-by-term integration of a series').

